\subsection{Avaliação fora do IMDb}
Aplicamos as seis CNNs do ensaio de tokenização, com os pesos congelados,
a 1.000 frases rotuladas da Amazon e 1.000 do Yelp, balanceadas por fonte
\cite{kotzias}. Os textos eram comentários curtos sobre produtos e serviços,
com medianas de 9 e 10 palavras, respectivamente.

\begin{tabularx}{\linewidth}{@{}Xrrrrrr@{}}
\toprule
 & \multicolumn{3}{c}{\tabelahead{CNN por palavras}}
 & \multicolumn{3}{c}{\tabelahead{CNN com BPE}}\\
\cmidrule(lr){2-4}\cmidrule(l){5-7}
\tabelahead{Fonte} & \tabelahead{Acurácia} & \tabelahead{F1 macro} & \tabelahead{CE}
 & \tabelahead{Acurácia} & \tabelahead{F1 macro} & \tabelahead{CE}\\
\midrule
\LinhasForaDominio
\bottomrule
\end{tabularx}
\legendinha{Tabela 10. Média das três sementes por tokenização.
\fonteatual{seminario/dados/fora_dominio_resultados.json}{Métricas por fonte e semente}.}

Em dois exemplos do YouTube sem rótulo de referência, o TF-IDF previu
positiva nos dois. As seis CNNs e o Qwen também leram o elogio seguido
de \textit{No one}, em \textit{The Room}, como positivo. No comentário
irônico de \textit{Suicide Squad}, cinco CNNs deram positiva e o Qwen
deu negativa nos dois modos de geração.
\fonteatual{seminario/dados/tfidf_youtube_resultados.json}{TF-IDF}.
\fonteatual{seminario/dados/fora_dominio_resultados.json}{CNN}.
\fonteatual{seminario/dados/qwen_youtube_resultados.json}{Qwen}.
